% Notes of the Montreal seminar, 1970 (cote 129, typed part): archive cote 129, site page 69 (PDF page 70).
% typescript. Draft compared: none.
% AI-transcribed by Claude Opus 5.5 (claude-opus-5-5, adaptive thinking, effort high), Message Batches API, batch mtl129_full (msgbatch_01AGnDKHwy6ATCM7ehUa6263), 2026-11-13; not yet checked by eye.
\scanpage{129/69}
\mspage{\hand{III-14}}

Un tore  T  est p-divisible et  T(1)  est un groupe fini localement libre
de rang  $p^{r}$,  r  étant la dimension relative de  T  sur  S  (ces proprié-%
tés sont locales et sont vérifiées par  $\mathbb{G}^{r}_{m,S}$).  On en déduit que  $T(\infty)$
est un groupe de  BT  de rang  r.

Comme exemple particulier important, on a le groupe de  BT
\[ \hand{\mu} = \mathbb{G}_{m,S}(\infty) = \varinjlim_{n} \hand{\mu}_{p^{n},S} \ . \]
\margin{$\times$ \ $\mu$}

6.4.  \uline{Exemple.}  Soit  G  un schéma en groupes sur  S,  extension d'un
schéma abélien  A  par un tore  T.  Alors  G  est p-divisible,  G(1)  est
fini localement libre sur  S  et  $G(\infty)$  est un groupe de  BT  de rang
2d\hand{+}r  (d  et  r  désignant respectivement les dimensions relatives de  A
et  T  sur  S).  De plus, on a une suite exacte
\[ 0 \to T(\infty) \to G(\infty) \to A(\infty) \to 0 \ . \]

6.5.  Comme y incite l'exemple (6.3), on dira qu'un groupe de  BT
G  est \uline{toroïdal} si  G(1)  est de type multiplicatif.  Il est équivalent
de dire que tous les  G(n)  sont de type multiplicatif, ou encore que le
dual  $G^{*}$  (5.5) est ind-étale.

Le foncteur  $G \mapsto T_{p}(G^{*})$  définit une anti-équivalence entre la
catégorie des groupes de  BT  toroidaux et celle des faisceaux p-adiques
constants tordus sans torsion.  Composant ce foncteur avec la dualité
dans les faisceaux p-adiques constants tordus sans torsion
\[ T_{p}(G^{*}) \mapsto \underline{\mathrm{Hom}}_{gr}(G^{*},\mathbb{Z}_{p}) \overset{\text{def}}{=} (\underline{\mathrm{Hom}}_{gr}(G(n)^{*},\Lambda_{n}))_{n\geq 1} \ , \]

\mspage{\hand{69}}
